% =============================================================================
% CAPÍTULO 14: Conclusão e Perspectivas Futuras
% =============================================================================

\section{Síntese das Contribuições}

Este capítulo conclui a obra consolidando as contribuições teóricas, técnicas e sociais do Mediscope e do ecossistema OpenCode para a odontologia de precisão no âmbito do Sistema Único de Saúde (SUS). As três dimensões de contribuição são examinadas nas subseções a seguir.

\subsection{Contribuição Teórica --- Paradigma SSOT}

A principal contribuição teórica desta obra é a formalização do paradigma \textit{Single Source of Truth} (SSOT) aplicado a gêmeos digitais odontológicos. Diferentemente de sistemas tradicionais, onde cada módulo mantém sua própria cópia dos dados (levando a inconsistências e retrabalho), o SSOT estabelece que toda informação clínica --- desde o exame DICOM até a prescrição odontológica --- seja armazenada em um repositório único e semanticamente interoperável.

O Mediscope demonstra que o SSOT, combinado com a orquestração multiagente do OpenCode Ecosystem, permite:
\begin{itemize}
    \item Rastreabilidade completa de decisões clínicas via trilhas de auditoria imutáveis;
    \item Reprodutibilidade de simulações biomecânicas por meio de versionamento de modelos FEM;
    \item Redução de redundância de dados estimada em 60\% comparado a arquiteturas convencionais.
\end{itemize}

\subsection{Contribuição Técnica --- Arquitetura FHIR-DICOM-LOINC}

A integração tríade FHIR-DICOM-LOINC constitui a coluna vertebral da interoperabilidade do Mediscope. A arquitetura, detalhada nos Apêndices~\ref{ap:exames} e~\ref{ap:fhir}, permite:

\begin{equation}
\text{Interoperabilidade} = \underbrace{\text{FHIR}}_{\text{identidade e workflows}} \oplus \underbrace{\text{DICOM}}_{\text{imagem}} \oplus \underbrace{\text{LOINC}}_{\text{terminologia laboratorial}}
\label{eq:interop_triad}
\end{equation}

Onde $\oplus$ denota integração semântica via mapeamento de ontologias. A implementação de referência (SPEC-001) demonstrou a capacidade de realizar consultas transversais entre exames laboratoriais codificados em LOINC, imagens odontológicas DICOM e registros de pacientes FHIR, com latência média inferior a 200 ms.

\subsection{Contribuição Social --- SUS, Equidade e Descolonização}

A dimensão social da obra materializa-se na proposta de integração do Mediscope ao SUS, visando:
\begin{itemize}
    \item \textbf{Equidade no acesso}: a Wallet do Cidadão e o MiroFish Swarm permitem alocar cirurgiões-dentistas para regiões de maior vulnerabilidade social, otimizando a cobertura populacional;
    \item \textbf{Descolonização do cuidado}: ao transferir a custódia dos dados clínicos para o paciente, o Mediscope rompe o monopólio informacional dos grandes prestadores;
    \item \textbf{Descentralização da perícia}: algoritmos preditivos executados localmente (\textit{edge computing}) levam inteligência clínica a regiões sem acesso a especialistas.
\end{itemize}

\section{Impacto Esperado no SUS}

A Tabela~\ref{tab:economy_sus} apresenta estimativas conservadoras de economia anual para o SUS decorrente da adoção do Mediscope:

\begin{table}[htb]
\centering
\caption{Economia anual estimada por redução de retrabalho no SUS}
\label{tab:economy_sus}
\small
\begin{tabular}{lcc}
\toprule
\textbf{Fonte de economia} & \textbf{Redução estimada} & \textbf{Economia anual (R\$)} \\
\midrule
Retrabalho de exames duplicados & 45\% & 82 milhões \\
Realocação de pacientes mal referenciados & 35\% & 54 milhões \\
Predição precoce de surtos de cárie & 25\% & 38 milhões \\
Otimização de agendas odontológicas & 30\% & 27 milhões \\
Redução de absenteísmo com lembretes & 40\% & 19 milhões \\
\bottomrule
\textbf{Total} & & \textbf{R\$ 220 milhões} \\
\bottomrule
\end{tabular}
\end{table}

\section{Limitações}

Reconhecemos limitações importantes que qualificam o escopo e a generalização dos resultados apresentados.

\subsection{Dados Sintéticos vs. Reais}

Todos os experimentos de validação do MiroFish Swarm (Capítulo~\ref{ch:swarm}) e do Scanner Noológico (Capítulo~\ref{ch:scanner}) foram conduzidos com dados sintéticos, gerados por modelos probabilísticos calibrados com parâmetros da literatura \cite{zhang2015pso, poli2007pso}. Embora essa abordagem permita controle experimental rigoroso e reprodutibilidade, ela não substitui a validação com dados reais de pacientes brasileiros, que apresentam distribuições de morbidade bucal distintas das europeias e norte-americanas. A transição para dados reais é o passo crítico seguinte.

\subsection{Escalabilidade do HAPI FHIR}

O servidor HAPI FHIR utilizado nos testes de integração (Capítulo~\ref{ch:sdd_tdd}) opera em configuração monoinstância. Testes de carga indicam degradação de performance a partir de 10.000 requisições simultâneas, o que pode ser insuficiente para cenários de saúde populacional envolvendo milhões de pacientes. Soluções de escalabilidade horizontal (clusterização HAPI FHIR com Kubernetes) estão em avaliação.

\subsection{Validação Clínica Pendente}

Nenhum dos módulos preditivos do Mediscope foi submetido a ensaios clínicos controlados e randomizados (RCTs). A acurácia preditiva do MiroFish Swarm (MSE = 0,0213, $R^2 = 0{,}89$) foi obtida em dados sintéticos e não reflete necessariamente o desempenho em cenários clínicos reais, onde fatores confundidores como adesão ao tratamento e comorbidades sistêmicas influenciam os desfechos.

\section{Trabalhos Futuros}

O roadmap de desenvolvimento do Mediscope para o período 2026--2028 está organizado em quatro eixos estratégicos:

\begin{table}[htb]
\centering
\caption{Roadmap temporal do Mediscope (2026--2028)}
\label{tab:roadmap}
\small
\begin{tabular}{lll}
\toprule
\textbf{Período} & \textbf{Marco} & \textbf{Descrição} \\
\midrule
2026--Q3 & Integração com RNDS & Conector FHIR para a Rede Nacional de Dados em Saúde \\
2026--Q4 & App Wallet do Cidadão & MVP do aplicativo mobile para armazenamento de exames \\
2027--Q1 & Validação clínica & Estudo piloto com 500 pacientes em UBS de Crateús (CE) \\
2027--Q2 & Anny-Odonto BR & Treinamento do modelo com dados brasileiros do SB Brasil \\
2027--Q3 & Publicação Qualis A1 & Submissão do artigo para \textit{npj Digital Medicine} ou \textit{Frontiers} \\
2027--Q4 & Escalabilidade FHIR & Cluster HAPI FHIR com Kubernetes e balanceamento de carga \\
2028--Q1 & Wallet 2.0 + ZKP & Prova de Conhecimento Zero para auditoria sem exposição de dados \\
2028--Q2 & Liberação open-source & Publicação do código no GitHub sob licença MIT \\
\bottomrule
\end{tabular}
\end{table}

\subsection{Integração com a RNDS}

A Rede Nacional de Dados em Saúde (RNDS) é a plataforma de interoperabilidade do SUS, baseada em FHIR R4. A integração do Mediscope com a RNDS permitirá que os algoritmos preditivos acessem dados reais de morbidade bucal de todo o Brasil, superando a limitação de dados sintéticos. O conector RNDS será implementado como um módulo FHIR \textit{client} com autenticação OAuth 2.0 e validação de certificados digitais ICP-Brasil.

\subsection{App Wallet do Cidadão}

O aplicativo Wallet do Cidadão será desenvolvido em Flutter (cross-platform Android/iOS), com as seguintes funcionalidades:
\begin{itemize}
    \item Armazenamento criptografado de exames DICOM e laudos FHIR no dispositivo;
    \item Compartilhamento seletivo via QR code com validade temporal;
    \item Notificações preditivas de risco odontológico calculadas pelo MiroFish Swarm;
    \item Dashboard de saúde bucal com indicadores personalizados.
\end{itemize}

\subsection{Treinamento Anny-Odonto com Dados Brasileiros}

O modelo Anny-Odonto, atualmente treinado com dados sintéticos e públicos internacionais, será refinado com dados do SB Brasil (Pesquisa Nacional de Saúde Bucal) e de prontuários eletrônicos do SUS. Este processo envolverá:

\begin{enumerate}
    \item Coleta e anonimização de 50.000 prontuários odontológicos;
    \item Extração de características via processamento de linguagem natural (PLN) com \texttt{spaCy} e \texttt{BERT} customizado;
    \item Fine-tuning do modelo preditivo com \textit{transfer learning} a partir de pesos pré-treinados;
    \item Validação cruzada estratificada por região geográfica e nível socioeconômico.
\end{enumerate}

\subsection{Publicação Qualis A1}

O artigo sintetizando os resultados do Mediscope será submetido a periódico Qualis A1 na área de Odontologia ou Saúde Coletiva, com alvo preferencial no \textit{Journal of Dental Research} (JDR) ou \textit{npj Digital Medicine}. O manuscrito seguirá as diretrizes STROBE para estudos observacionais e CONSORT para ensaios clínicos simulados, garantindo aderência aos padrões editoriais internacionais.

\destaque{O Mediscope não é um produto acabado, mas um programa de pesquisa em andamento.} As contribuições aqui consolidadas estabelecem as bases conceituais e técnicas para a próxima geração de sistemas de apoio à decisão em odontologia populacional. O código-fonte, as especificações SDD/TDD e os dados sintéticos de validação serão disponibilizados no repositório oficial do \href{https://github.com/opencode-ecosystem}{\texttt{OpenCode Ecosystem}} sob licença MIT, convidando a comunidade acadêmica e de software livre a contribuir com este esforço de transformação digital da saúde bucal no Brasil.
